\chapter{Introduction}
\label{ch:introduction}

% Research question — frame the thesis as a preliminary answer to a concrete question.
% Opening: the threat, the gap in defenses, and why AMR is a candidate solution.

\section{Prompt Injection in LLM Agents}
\label{sec:intro-threat}

% Describe the threat: LLM agents that consume external tool outputs are vulnerable
% to injected instructions embedded in that data.  Direct and indirect injection.
% Concrete example: banking agent reads a landlord notice containing a hidden payment command.

\section{Limitations of Existing Defenses}
\label{sec:intro-defenses}

% Keyword filters: surface-level, bypass-trivial.
% LLM judges: operate on raw text, themselves injectable, no formal guarantees.
% Gap: no semantically-grounded, injection-safe detection with formal motivation.

\section{Research Question}
\label{sec:intro-rq}

% Explicit question: Can AMR-based static analysis help detect prompt injection attacks?
% Why AMR: lossy compression as a security feature — synonyms, paraphrases, and
% obfuscation collapse to the same predicate-argument graph.
% Scope: we answer this question in a narrow but meaningful context (AgentDojo
% banking benchmark with explicit system policies).

\section{Contributions}
\label{sec:intro-contributions}

% Bullet list:
% 1. Scoped formal definition of prompt injection as AMR-graph contradiction.
% 2. A static firewall with four detection rules and a minimal LLM cost model.
% 3. An authorization annotation mechanism that handles the data/instruction boundary.
% 4. Empirical evaluation on a custom AgentDojo attack suite (X attacks, Y benign tasks).
% 5. Analysis of structural graph similarity (smatch) and why it fails as a primary rule.

\section{Thesis Outline}
\label{sec:intro-outline}

% One paragraph mapping chapters to contributions.
